\documentclass[11pt]{article}
\usepackage[a4paper,margin=0.75in]{geometry}
\usepackage{amsmath,amssymb,amsthm,hyperref}
\hypersetup{hidelinks}
\newtheorem{theorem}{Theorem}
\newcommand{\M}{\mathrm{M}}
\title{A counterexample to a squared operator P\'olya--Szeg\H{o} conjecture}
\author{Research draft: independent review pending}
\date{October 27, 2026}
\begin{document}
\maketitle
\begin{abstract}
We give positive definite matrices in $\M_4(\mathbb C)$ and a unital positive
linear map to $\M_2(\mathbb C)$ that violate the first inequality of
Moslehian--Fu's Conjecture~2.16. The inputs have spectra in $[1,1024]$.
Two explicit geometric-mean formulas reduce the obstruction to a negative
quadratic form certified by rational bounds on two square roots.
The finite-dimensional refutation has a closed Lean certificate.
\end{abstract}
\noindent\textbf{Draft status.} Discovery credit is zero pending independent
source, mathematics, formal-scope, novelty and manuscript review, followed by
the final extracted replay. No priority claim is made.

For positive definite matrices, write
$X\sharp Y=X^{1/2}(X^{-1/2}YX^{-1/2})^{1/2}X^{1/2}$.
The first display of \cite[Conjecture~2.16, p.~9]{MF} proposes
\begin{equation}\label{claim}
 (\Phi(U)\sharp\Phi(V))^2\leq
 \frac14\left(\sqrt{\frac{M_1M_2}{m_1m_2}}+
 \sqrt{\frac{m_1m_2}{M_1M_2}}\right)^2\Phi(U\sharp V)^2
\end{equation}
for every unital positive linear map $\Phi$, where
$m_1^2I\leq U\leq M_1^2I$ and $m_2^2I\leq V\leq M_2^2I$,
with $0<m_i\leq M_i$. The order is the positive-semidefinite order.

\begin{theorem}\label{counter}
Inequality~\eqref{claim} fails for $m_1=m_2=1$, $M_1=M_2=32$,
$U,V\in\M_4(\mathbb C)$ and a unital positive complex-linear
$\Phi:\M_4(\mathbb C)\to\M_2(\mathbb C)$.
\end{theorem}
\begin{proof}
Let $I=I_2$ and put
\[
 A=\frac12\begin{pmatrix}1025&-1023\\-1023&1025\end{pmatrix},\quad
 B=\begin{pmatrix}1024&0\\0&1\end{pmatrix},\quad
 U=\operatorname{diag}(A,B),\quad V=\operatorname{diag}(B,I).
\]
Both $A$ and $B$ have eigenvalues $1,1024$, so the required input bounds hold.
For an arbitrary $2\times2$ block matrix $Z$, define
$\Phi(Z)=(Z_{11}+Z_{22})/2$. This is the average of two compression maps,
hence is unital, complex-linear and positive (indeed completely positive).
Set
\[
 F=\Phi(U)=\frac14\begin{pmatrix}3073&-1023\\-1023&1027\end{pmatrix},\quad
 G=\Phi(V)=\begin{pmatrix}1025/2&0\\0&1\end{pmatrix},\quad
 T=\begin{pmatrix}32&0\\0&1\end{pmatrix}.
\]
Write
\begin{equation}\label{radicals}
 K=1054721,\quad a=\sqrt{2048/K},\quad b=\sqrt{4100/K},\quad
 r=1058821/K,\quad \rho=r+2b.
\end{equation}
The two geometric means needed below are
\begin{equation}\label{means}
 C:=\Phi(U\sharp V)=aF+T/2,\qquad
 D:=\Phi(U)\sharp\Phi(V)=\frac{G+bF}{\sqrt\rho}.
\end{equation}
To verify these as genuine positive geometric means, recall that $X\sharp Y$
is the unique positive definite solution $Z$ of $ZX^{-1}Z=Y$.
Direct rational multiplication gives
\begin{align*}
 (A+B)A^{-1}(A+B)&=(K/2048)B,& T B^{-1}T&=I,\\
 GF^{-1}G&=rG-(4100/K)F.
\end{align*}
Consequently $A\sharp B=a(A+B)$, $B\sharp I=T$, and
$(G+bF)F^{-1}(G+bF)=\rho G$. Every proposed solution is positive
definite. Applying the same characterization to the block-diagonal inputs
proves both identities in~\eqref{means}.

The coefficient in~\eqref{claim} is $\alpha^2$, where $\alpha=1025/64$.
For $v=(1,2)^\mathsf T$, exact multiplication yields
\begin{align}
 v^\mathsf T C^2v&=\frac{2048}{K}\frac{1058845}{8}
                   +a\frac{17463}{2}+257,\label{cform}\\
 v^\mathsf T D^2v&=
 \frac{\frac{1050641}{4}+b\frac{1056799}{4}
                  +\frac{4100}{K}\frac{1058845}{8}}{r+2b}.\label{dform}
\end{align}
Squaring positive rationals establishes
\[
 0\leq a\leq\frac{441}{10000},\qquad
 \frac{623}{10000}\leq b\leq\frac{624}{10000}.
\]
Use the upper bound for $a$ in~\eqref{cform}, the lower bound for $b$
in the numerator of~\eqref{dform}, and its upper bound in the denominator.
All the latter quantities are positive. Ordinary rational arithmetic then gives
\begin{equation}\label{gap}
 v^\mathsf T(\alpha^2 C^2-D^2)v
 \leq-\frac{1763270875137191128175507}{102858162597360173056}
 <-17000.
\end{equation}
Thus $\alpha^2C^2-D^2$ is not positive semidefinite, proving the theorem.
\end{proof}

The construction also fits an endomorphism convention for $\Phi$:
replace it by $Z\mapsto\operatorname{diag}(\Phi(Z),\Phi(Z))$ and test on
$(v,0)$. The two geometric means and their differences double block diagonally.

\section*{Formal scope and literature qualifications}
The certificate uses Lean~4.34.1 and mathlib revision
\texttt{d13f23b723b8a846827a245b89c10fc7d3f11612}.
Its $100$ declarations in $9$ modules passed a fresh isolated replay with
warnings treated as errors; only \texttt{propext}, \texttt{Classical.choice}
and \texttt{Quot.sound} occur. The endpoint
\texttt{source\_counterexample\_}\allowbreak\texttt{329636\_first} is a closed negation of the
universal finite-dimensional version of~\eqref{claim}. The positive map,
input bounds, functional-calculus square roots, Riccati identifications,
quadratic-form interpretation and strict sign are proved in Lean.
This finite instance suffices to refute the source's unrestricted operator
assertion. No optimal constant or additional discovery is claimed.

The second, Diaz--Metcalf, clause of the source conjecture is outside this
formal theorem. Related claimed counterexamples occur in \cite[Remark~1]{H}
and \cite[Remark~2.4]{E}; checking them requires their stated spectral bounds.
For example, at the bounds $m=1/50$, $M=11$ printed in \cite{H}, the determinant
of its claimed negative difference is actually
\[
 \frac{88033691903}{87120000}+\frac{10018833\sqrt6}{12800}>0.
\]
Its first leading minor is positive too. Our supplementary exact checks record
this discrepancy, without awarding credit for correcting it or assuming it
settles the broader literature. The present proposed outcome is Theorem~\ref{counter}.

The comparison uses the actual arXiv statement of \cite{MF}; the final
published LAA body remains unread because of access limits. The bounded prior-art
screen is not a novelty certificate: a close 2018 citing paper remains unread,
and independent review is pending. There is no claim of worldwide priority.

\small
\begin{thebibliography}{9}
\bibitem{MF} M.~S. Moslehian and X.~Fu,
\emph{Squaring operator P\'olya--Szeg\H{o} and Diaz--Metcalf type inequalities},
Linear Algebra Appl. 491 (2016), 73--82;
\href{https://arxiv.org/abs/1501.02939v1}{arXiv:1501.02939v1}.
\bibitem{H} D.~T. Hoa, D.~T.~H. Binh and T.~M. Ho,
\emph{On some matrix mean inequalities with Kantorovich constant},
Sci. Math. Jpn. 80 (2017), 139--151; online version received 2014, revised 2015.
\bibitem{E} D.~T. Hoa, M.~S. Moslehian, C.~Conde and P.~Zhang,
\emph{An extension of the P\'olya--Szeg\H{o} operator inequality},
\href{https://arxiv.org/abs/1610.04162v1}{arXiv:1610.04162v1} (2016).
\end{thebibliography}
\end{document}
